\problemname{The Siege}

\setlength{\parindent}{0pt}
\setlength{\parskip}{6pt}

You are Warchief Grukk, and today is a good day for war.
Your Horde stands ready to crush the Dwarven Underkingdom,
$N$~strongholds carved deep into the mountains, each sitting on rich
veins of gold.

Your goblin scouts (small, sneaky, and not terribly bright) have
somehow wormed their way into every stronghold through cracks and drain
pipes the dwarves forgot to seal. They came back clutching crumpled
maps scratched onto hide scraps, giggling about ``shiny rocks.''
The maps are crude but surprisingly accurate: each one shows the layout
of tunnels inside a stronghold.

Deep inside each stronghold sits the \emph{Gold Vein} (room~$0$),
where the richest ore waits to be mined. Once you conquer a stronghold,
your goblin miners will dig out the gold and load it into mine carts
that run through the tunnel network to the stronghold entrance
(room~$V_i - 1$), where your wolf-riders wait with wagons.

Each tunnel has a \emph{weight limit}: the maximum kilograms of gold
ore that a mine cart can carry through it per hour. Some tunnels are
wide enough for a troll to stroll through; others are barely
goblin-sized. Your miners will haul as much gold as physically possible
from the vein to the entrance, using every tunnel to its limit
simultaneously. You need to figure out how much gold each stronghold
can yield.

But strongholds do not fall for free. The dwarves fight like cornered
badgers behind their stone barricades. Each assault costs orc warriors
their lives. You can muster exactly $B$~warriors for the entire
campaign, and you must conquer your chosen strongholds before the elven
relief army arrives. Choose which strongholds to take to maximise the
total gold your miners extract. Each stronghold can only be taken once.

Present your \emph{Siege Plans} to the war council: the list of
strongholds the Horde will conquer.

\begin{Input}
The first line contains two integers $N$ and
$B$~$(1 \le N \le 50,\; 1 \le B \le 1000)$,
the number of dwarven strongholds and the number of orc warriors
available.

Then $N$ stronghold descriptions follow. Each stronghold starts with
a line containing:
\begin{itemize}
  \item a string $s_i$, the name of the stronghold
        (lowercase \texttt{a}--\texttt{z}, no whitespace,
        length $\le 20$),
  \item an integer $c_i$~$(1 \le c_i \le B)$, the number of orc
        warriors lost conquering it,
  \item an integer $V_i$~$(2 \le V_i \le 15)$, the number of rooms,
  \item an integer $E_i$~$(1 \le E_i \le 40)$, the number of tunnels.
\end{itemize}

The next $E_i$ lines each describe one tunnel with three integers
$u$,~$v$,~$w$~$(0 \le u, v < V_i,\; u \ne v,\; 1 \le w \le 100)$,
meaning a one-way tunnel from room~$u$ to room~$v$ with weight
limit~$w$.

Room~$0$ is the Gold Vein (source) and room~$V_i - 1$ is the
entrance (sink).
\end{Input}

\begin{Output}
Output the strongholds you should conquer to maximise total gold
extracted without exceeding the warrior budget~$B$.
Print one line per selected stronghold in the format:

\texttt{\quad <index> <name>}

\noindent
where \texttt{<index>} is the 0-based position of the stronghold in
the input. Lines must appear in ascending order of index.
It is guaranteed that the optimal subset is unique.
\end{Output}
